\section{晚清改革证据链事件锚点}
以下锚点连接灾荒、边疆、学校、公司、法律和宪政的补充记录，明确制度文本、公共信息和地方实施之间仍有需核验的距离。
\subsection{1870—1879}
\eventanchor{EQ-1875-GUANGXU-ACCESSION-DOCUMENTS}{围绕「光绪帝即位，慈禧、慈安继…}光绪元年（1875年），围绕「光绪帝即位，慈禧、慈安继续掌握关键政治影响」的上谕、奏折与部议在中央—地方行政链中流转。核心取证：《清德宗实录》光绪元年相关条；宫中朱批奏折、内阁大库题本与《清史稿》相关志传。这一锚点记录的是文书链：它能证明呈报、上谕、部议或照会进入机构流程，不能由此推出实际执行。来源保留：此行仅确认相关文书链和可追查的行政动作；不得由其直接推断政策有效、地方服从、民众态度或单一因果。
\eventanchor{EQ-1875-YUNNAN-MUSLIM-WAR-DOCUMENTS}{清廷围绕「杜文秀政权覆亡前后，…}光绪元年（1875年），清廷围绕「杜文秀政权覆亡前后，云南回民战争进入惨烈收束阶段」调遣军队、筹措饷械并处理战报，中央与地方的军政文书形成连续记录。核心取证：《清德宗实录》光绪元年相关条；军机处档、战事方略及地方奏报。这一锚点记录的是文书链：它能证明呈报、上谕、部议或照会进入机构流程，不能由此推出实际执行。来源保留：此行仅确认相关文书链和可追查的行政动作；不得由其直接推断政策有效、地方服从、民众态度或单一因果。
\eventanchor{EQ-1876-CHEFOO-CONVENTION-DOCUMENTS}{围绕「《烟台条约》签订，扩大英…}光绪二年（1876年），围绕「《烟台条约》签订，扩大英国在华通商和内地活动权利」的照会、条约文本、换约或地方执行报告分别进入外交文书链。核心取证：《清德宗实录》光绪二年相关条；《筹办夷务始末》相关年卷与中外照会、条约文本。这一锚点记录的是文书链：它能证明呈报、上谕、部议或照会进入机构流程，不能由此推出实际执行。来源保留：此行仅确认相关文书链和可追查的行政动作；不得由其直接推断政策有效、地方服从、民众态度或单一因果。
\eventanchor{EQ-1876-NORTH-CHINA-FAMINE-DOCUMENTS}{围绕「丁戊奇荒在华北及西北造成…}光绪二年（1876年），围绕「丁戊奇荒在华北及西北造成广泛饥荒、迁徙和救济行动」的地方呈报、案件或赈务记录将社会反应带入官府文书链。核心取证：《清德宗实录》光绪二年相关条；地方志、督抚奏折及地方档案。这一锚点记录的是文书链：它能证明呈报、上谕、部议或照会进入机构流程，不能由此推出实际执行。来源保留：此行仅确认相关文书链和可追查的行政动作；不得由其直接推断政策有效、地方服从、民众态度或单一因果。
\eventanchor{EQ-1877-FAMINE-RELIEF-NETWORKS-DOCUMENTS}{围绕「赈济网络在灾区展开，官办…}光绪三年（1877年），围绕「赈济网络在灾区展开，官办、绅办和传教团体的救助并存」的地方呈报、案件或赈务记录将社会反应带入官府文书链。核心取证：《清德宗实录》光绪三年相关条；地方志、督抚奏折及地方档案。这一锚点记录的是文书链：它能证明呈报、上谕、部议或照会进入机构流程，不能由此推出实际执行。来源保留：此行仅确认相关文书链和可追查的行政动作；不得由其直接推断政策有效、地方服从、民众态度或单一因果。
\eventanchor{EQ-1878-Zuo-ZONGTANG-XINJIANG-DOCUMENTS}{清廷围绕「左宗棠军收复新疆大部…}光绪四年（1878年），清廷围绕「左宗棠军收复新疆大部，阿古柏政权崩解」调遣军队、筹措饷械并处理战报，中央与地方的军政文书形成连续记录。核心取证：《清德宗实录》光绪四年相关条；军机处档、战事方略及地方奏报。这一锚点记录的是文书链：它能证明呈报、上谕、部议或照会进入机构流程，不能由此推出实际执行。来源保留：此行仅确认相关文书链和可追查的行政动作；不得由其直接推断政策有效、地方服从、民众态度或单一因果。
\eventanchor{EQ-1878-XINJIANG-PROVINCE-DEBATE-DOCUMENTS}{围绕「朝廷讨论新疆设省、军政与…}光绪四年（1878年），围绕「朝廷讨论新疆设省、军政与财政安排」的上谕、奏折与部议在中央—地方行政链中流转。核心取证：《清德宗实录》光绪四年相关条；宫中朱批奏折、内阁大库题本与《清史稿》相关志传。这一锚点记录的是文书链：它能证明呈报、上谕、部议或照会进入机构流程，不能由此推出实际执行。来源保留：此行仅确认相关文书链和可追查的行政动作；不得由其直接推断政策有效、地方服从、民众态度或单一因果。
\eventanchor{EQ-1879-RYUKYU-ANNEXATION-DOCUMENTS}{围绕「日本吞并琉球，清廷的交涉…}光绪五年（1879年），围绕「日本吞并琉球，清廷的交涉未能恢复旧有册封关系」的照会、条约文本、换约或地方执行报告分别进入外交文书链。核心取证：《清德宗实录》光绪五年相关条；《筹办夷务始末》相关年卷与中外照会、条约文本。这一锚点记录的是文书链：它能证明呈报、上谕、部议或照会进入机构流程，不能由此推出实际执行。来源保留：此行仅确认相关文书链和可追查的行政动作；不得由其直接推断政策有效、地方服从、民众态度或单一因果。
\subsection{1880—1889}
\eventanchor{EQ-1880-ILI-NEGOTIATIONS-DOCUMENTS}{围绕「清俄围绕伊犁归还和边界条…}光绪六年（1880年），围绕「清俄围绕伊犁归还和边界条件持续谈判」的照会、条约文本、换约或地方执行报告分别进入外交文书链。核心取证：《清德宗实录》光绪六年相关条；《筹办夷务始末》相关年卷与中外照会、条约文本。这一锚点记录的是文书链：它能证明呈报、上谕、部议或照会进入机构流程，不能由此推出实际执行。来源保留：此行仅确认相关文书链和可追查的行政动作；不得由其直接推断政策有效、地方服从、民众态度或单一因果。
\eventanchor{EQ-1881-SAINT-PETERSBURG-TREATY-DOCUMENTS}{围绕「《中俄伊犁条约》签订，伊…}光绪七年（1881年），围绕「《中俄伊犁条约》签订，伊犁大部归还清廷并附带领土、赔款和贸易条件」的照会、条约文本、换约或地方执行报告分别进入外交文书链。核心取证：《清德宗实录》光绪七年相关条；《筹办夷务始末》相关年卷与中外照会、条约文本。这一锚点记录的是文书链：它能证明呈报、上谕、部议或照会进入机构流程，不能由此推出实际执行。来源保留：此行仅确认相关文书链和可追查的行政动作；不得由其直接推断政策有效、地方服从、民众态度或单一因果。
\eventanchor{EQ-1882-IMO-INCIDENT-DOCUMENTS}{围绕「朝鲜壬午兵变后清军入朝，…}光绪八年（1882年），围绕「朝鲜壬午兵变后清军入朝，日本也扩大影响」的照会、条约文本、换约或地方执行报告分别进入外交文书链。核心取证：《清德宗实录》光绪八年相关条；《筹办夷务始末》相关年卷与中外照会、条约文本。这一锚点记录的是文书链：它能证明呈报、上谕、部议或照会进入机构流程，不能由此推出实际执行。来源保留：此行仅确认相关文书链和可追查的行政动作；不得由其直接推断政策有效、地方服从、民众态度或单一因果。
\eventanchor{EQ-1883-SINO-FRENCH-WAR-BEGIN-DOCUMENTS}{清廷围绕「中法因越南北部和宗藩…}光绪九年（1883年），清廷围绕「中法因越南北部和宗藩关系冲突升级为战争」调遣军队、筹措饷械并处理战报，中央与地方的军政文书形成连续记录。核心取证：《清德宗实录》光绪九年相关条；军机处档、战事方略及地方奏报。这一锚点记录的是文书链：它能证明呈报、上谕、部议或照会进入机构流程，不能由此推出实际执行。来源保留：此行仅确认相关文书链和可追查的行政动作；不得由其直接推断政策有效、地方服从、民众态度或单一因果。
\eventanchor{EQ-1884-FUZHOU-NAVAL-BATTLE-DOCUMENTS}{清廷围绕「马江海战中福建水师遭…}光绪十年（1884年），清廷围绕「马江海战中福建水师遭法国舰队重创」调遣军队、筹措饷械并处理战报，中央与地方的军政文书形成连续记录。核心取证：《清德宗实录》光绪十年相关条；军机处档、战事方略及地方奏报。这一锚点记录的是文书链：它能证明呈报、上谕、部议或照会进入机构流程，不能由此推出实际执行。来源保留：此行仅确认相关文书链和可追查的行政动作；不得由其直接推断政策有效、地方服从、民众态度或单一因果。
\eventanchor{EQ-1884-TAIWAN-BLOCKADE-DOCUMENTS}{清廷围绕「法军进攻台湾北部并封…}光绪十年（1884年），清廷围绕「法军进攻台湾北部并封锁，台湾成为中法战争重要战场」调遣军队、筹措饷械并处理战报，中央与地方的军政文书形成连续记录。核心取证：《清德宗实录》光绪十年相关条；军机处档、战事方略及地方奏报。这一锚点记录的是文书链：它能证明呈报、上谕、部议或照会进入机构流程，不能由此推出实际执行。来源保留：此行仅确认相关文书链和可追查的行政动作；不得由其直接推断政策有效、地方服从、民众态度或单一因果。
\eventanchor{EQ-1885-TIANJIN-CHINA-FRANCE-TREATY-DOCUMENTS}{围绕「《中法新约》确认越南脱离…}光绪十一年（1885年），围绕「《中法新约》确认越南脱离清朝宗藩体系，战争结束」的照会、条约文本、换约或地方执行报告分别进入外交文书链。核心取证：《清德宗实录》光绪十一年相关条；《筹办夷务始末》相关年卷与中外照会、条约文本。这一锚点记录的是文书链：它能证明呈报、上谕、部议或照会进入机构流程，不能由此推出实际执行。来源保留：此行仅确认相关文书链和可追查的行政动作；不得由其直接推断政策有效、地方服从、民众态度或单一因果。
\eventanchor{EQ-1885-TAIWAN-PROVINCE-DOCUMENTS}{围绕「清廷决定将台湾建省，加强…}光绪十一年（1885年），围绕「清廷决定将台湾建省，加强行政、军防和基础设施经营」的上谕、奏折与部议在中央—地方行政链中流转。核心取证：《清德宗实录》光绪十一年相关条；宫中朱批奏折、内阁大库题本与《清史稿》相关志传。这一锚点记录的是文书链：它能证明呈报、上谕、部议或照会进入机构流程，不能由此推出实际执行。来源保留：此行仅确认相关文书链和可追查的行政动作；不得由其直接推断政策有效、地方服从、民众态度或单一因果。
\eventanchor{EQ-1886-BEIYANG-NAVY-EXPANSION-DOCUMENTS}{清廷围绕「北洋海军和沿海防务建…}光绪十二年（1886年），清廷围绕「北洋海军和沿海防务建设持续推进」调遣军队、筹措饷械并处理战报，中央与地方的军政文书形成连续记录。核心取证：《清德宗实录》光绪十二年相关条；军机处档、战事方略及地方奏报。这一锚点记录的是文书链：它能证明呈报、上谕、部议或照会进入机构流程，不能由此推出实际执行。来源保留：此行仅确认相关文书链和可追查的行政动作；不得由其直接推断政策有效、地方服从、民众态度或单一因果。
\eventanchor{EQ-1887-TAIWAN-RAILWAY-PLANNING-DOCUMENTS}{围绕「台湾巡抚刘铭传推动铁路、…}光绪十三年（1887年），围绕「台湾巡抚刘铭传推动铁路、电报、矿务和新式行政项目」的经费、税收、赈济或工程呈报经由户部与地方衙门处理。核心取证：《清德宗实录》光绪十三年相关条；户部奏销、盐政、河工或海关档案。这一锚点记录的是文书链：它能证明呈报、上谕、部议或照会进入机构流程，不能由此推出实际执行。来源保留：此行仅确认相关文书链和可追查的行政动作；不得由其直接推断政策有效、地方服从、民众态度或单一因果。
\eventanchor{EQ-1888-NAVAL-AND-ARMY-REFORM-DOCUMENTS}{围绕「清廷讨论海军、陆军训练和…}光绪十四年（1888年），围绕「清廷讨论海军、陆军训练和军工体系的经费与制度问题」的上谕、奏折与部议在中央—地方行政链中流转。核心取证：《清德宗实录》光绪十四年相关条；宫中朱批奏折、内阁大库题本与《清史稿》相关志传。这一锚点记录的是文书链：它能证明呈报、上谕、部议或照会进入机构流程，不能由此推出实际执行。来源保留：此行仅确认相关文书链和可追查的行政动作；不得由其直接推断政策有效、地方服从、民众态度或单一因果。
\eventanchor{EQ-1889-GUANGXU-MARRIAGE-DOCUMENTS}{围绕「光绪帝大婚并开始形式上的…}光绪十五年（1889年），围绕「光绪帝大婚并开始形式上的亲政安排」的上谕、奏折与部议在中央—地方行政链中流转。核心取证：《清德宗实录》光绪十五年相关条；宫中朱批奏折、内阁大库题本与《清史稿》相关志传。这一锚点记录的是文书链：它能证明呈报、上谕、部议或照会进入机构流程，不能由此推出实际执行。来源保留：此行仅确认相关文书链和可追查的行政动作；不得由其直接推断政策有效、地方服从、民众态度或单一因果。
\subsection{1890—1899}
\eventanchor{EQ-1890-KOREAN-POLITICAL-CONFLICT-DOCUMENTS}{围绕「朝鲜内政纷争持续牵动清日…}光绪十六年（1890年），围绕「朝鲜内政纷争持续牵动清日两国在半岛的影响力竞争」的照会、条约文本、换约或地方执行报告分别进入外交文书链。核心取证：《清德宗实录》光绪十六年相关条；《筹办夷务始末》相关年卷与中外照会、条约文本。这一锚点记录的是文书链：它能证明呈报、上谕、部议或照会进入机构流程，不能由此推出实际执行。来源保留：此行仅确认相关文书链和可追查的行政动作；不得由其直接推断政策有效、地方服从、民众态度或单一因果。
\eventanchor{EQ-1891-MISSIONARY-CONFLICTS-DOCUMENTS}{围绕「各地教案和宗教纠纷继续引…}光绪十七年（1891年），围绕「各地教案和宗教纠纷继续引发地方治理与外交压力」的地方呈报、案件或赈务记录将社会反应带入官府文书链。核心取证：《清德宗实录》光绪十七年相关条；地方志、督抚奏折及地方档案。这一锚点记录的是文书链：它能证明呈报、上谕、部议或照会进入机构流程，不能由此推出实际执行。来源保留：此行仅确认相关文书链和可追查的行政动作；不得由其直接推断政策有效、地方服从、民众态度或单一因果。
\eventanchor{EQ-1892-RAILWAY-DEBATE-DOCUMENTS}{围绕「官员、商人和士绅围绕铁路…}光绪十八年（1892年），围绕「官员、商人和士绅围绕铁路修筑、主权与资金展开争论」的经费、税收、赈济或工程呈报经由户部与地方衙门处理。核心取证：《清德宗实录》光绪十八年相关条；户部奏销、盐政、河工或海关档案。这一锚点记录的是文书链：它能证明呈报、上谕、部议或照会进入机构流程，不能由此推出实际执行。来源保留：此行仅确认相关文书链和可追查的行政动作；不得由其直接推断政策有效、地方服从、民众态度或单一因果。
\eventanchor{EQ-1893-KOREAN-REFORM-CRISIS-DOCUMENTS}{围绕「朝鲜改革与政治危机使清日…}光绪十九年（1893年），围绕「朝鲜改革与政治危机使清日关系持续紧张」的照会、条约文本、换约或地方执行报告分别进入外交文书链。核心取证：《清德宗实录》光绪十九年相关条；《筹办夷务始末》相关年卷与中外照会、条约文本。这一锚点记录的是文书链：它能证明呈报、上谕、部议或照会进入机构流程，不能由此推出实际执行。来源保留：此行仅确认相关文书链和可追查的行政动作；不得由其直接推断政策有效、地方服从、民众态度或单一因果。
\eventanchor{EQ-1894-DONGHAK-PEASANT-WAR-DOCUMENTS}{清廷围绕「东学农民战争促使清日…}光绪二十年（1894年），清廷围绕「东学农民战争促使清日两国派兵入朝，危机迅速国际化」调遣军队、筹措饷械并处理战报，中央与地方的军政文书形成连续记录。核心取证：《清德宗实录》光绪二十年相关条；军机处档、战事方略及地方奏报。这一锚点记录的是文书链：它能证明呈报、上谕、部议或照会进入机构流程，不能由此推出实际执行。来源保留：此行仅确认相关文书链和可追查的行政动作；不得由其直接推断政策有效、地方服从、民众态度或单一因果。
\eventanchor{EQ-1894-FIRST-SINO-JAPANESE-WAR-DOCUMENTS}{清廷围绕「甲午战争爆发，清日围…}光绪二十年（1894年），清廷围绕「甲午战争爆发，清日围绕朝鲜和区域秩序展开海陆战」调遣军队、筹措饷械并处理战报，中央与地方的军政文书形成连续记录。核心取证：《清德宗实录》光绪二十年相关条；军机处档、战事方略及地方奏报。这一锚点记录的是文书链：它能证明呈报、上谕、部议或照会进入机构流程，不能由此推出实际执行。来源保留：此行仅确认相关文书链和可追查的行政动作；不得由其直接推断政策有效、地方服从、民众态度或单一因果。
\eventanchor{EQ-1894-YALU-NAVAL-BATTLE-DOCUMENTS}{清廷围绕「黄海海战后北洋舰队受…}光绪二十年（1894年），清廷围绕「黄海海战后北洋舰队受创，海军力量对比进一步恶化」调遣军队、筹措饷械并处理战报，中央与地方的军政文书形成连续记录。核心取证：《清德宗实录》光绪二十年相关条；军机处档、战事方略及地方奏报。这一锚点记录的是文书链：它能证明呈报、上谕、部议或照会进入机构流程，不能由此推出实际执行。来源保留：此行仅确认相关文书链和可追查的行政动作；不得由其直接推断政策有效、地方服从、民众态度或单一因果。
\eventanchor{EQ-1895-SHIMONOSEKI-TREATY-DOCUMENTS}{围绕「《马关条约》签订，清割让…}光绪二十一年（1895年），围绕「《马关条约》签订，清割让台湾、澎湖并承认朝鲜独立，赔款与通商条件加重」的照会、条约文本、换约或地方执行报告分别进入外交文书链。核心取证：《清德宗实录》光绪二十一年相关条；《筹办夷务始末》相关年卷与中外照会、条约文本。这一锚点记录的是文书链：它能证明呈报、上谕、部议或照会进入机构流程，不能由此推出实际执行。来源保留：此行仅确认相关文书链和可追查的行政动作；不得由其直接推断政策有效、地方服从、民众态度或单一因果。
\eventanchor{EQ-1895-TAIWAN-REPUBLIC-RESISTANCE-DOCUMENTS}{清廷围绕「台湾割让后地方力量建…}光绪二十一年（1895年），清廷围绕「台湾割让后地方力量建立台湾民主国并抵抗日军登陆」调遣军队、筹措饷械并处理战报，中央与地方的军政文书形成连续记录。核心取证：《清德宗实录》光绪二十一年相关条；军机处档、战事方略及地方奏报。这一锚点记录的是文书链：它能证明呈报、上谕、部议或照会进入机构流程，不能由此推出实际执行。来源保留：此行仅确认相关文书链和可追查的行政动作；不得由其直接推断政策有效、地方服从、民众态度或单一因果。
\eventanchor{EQ-1895-TRIPLE-INTERVENTION-DOCUMENTS}{围绕「三国干涉还辽改变战后列强…}光绪二十一年（1895年），围绕「三国干涉还辽改变战后列强竞争和清廷外交处境」的照会、条约文本、换约或地方执行报告分别进入外交文书链。核心取证：《清德宗实录》光绪二十一年相关条；《筹办夷务始末》相关年卷与中外照会、条约文本。这一锚点记录的是文书链：它能证明呈报、上谕、部议或照会进入机构流程，不能由此推出实际执行。来源保留：此行仅确认相关文书链和可追查的行政动作；不得由其直接推断政策有效、地方服从、民众态度或单一因果。
\eventanchor{EQ-1896-LI-LOBANOV-TREATY-DOCUMENTS}{围绕「《中俄密约》及东清铁路安…}光绪二十二年（1896年），围绕「《中俄密约》及东清铁路安排加强俄国在东北的影响」的照会、条约文本、换约或地方执行报告分别进入外交文书链。核心取证：《清德宗实录》光绪二十二年相关条；《筹办夷务始末》相关年卷与中外照会、条约文本。这一锚点记录的是文书链：它能证明呈报、上谕、部议或照会进入机构流程，不能由此推出实际执行。来源保留：此行仅确认相关文书链和可追查的行政动作；不得由其直接推断政策有效、地方服从、民众态度或单一因果。
\eventanchor{EQ-1896-MODERN-EDUCATION-EXPANSION-DOCUMENTS}{围绕「甲午战后新式学堂、译书与…}光绪二十二年（1896年），围绕「甲午战后新式学堂、译书与留学讨论加速」的禁令、编纂、教育或传播活动留下中央和地方文本记录。核心取证：《清德宗实录》光绪二十二年相关条；内阁档、文集或报刊相关记录。这一锚点记录的是文书链：它能证明呈报、上谕、部议或照会进入机构流程，不能由此推出实际执行。来源保留：此行仅确认相关文书链和可追查的行政动作；不得由其直接推断政策有效、地方服从、民众态度或单一因果。
\eventanchor{EQ-1897-JIAOZHOU-BAY-OCCUPATION-DOCUMENTS}{围绕「德国以巨野教案为由占领胶…}光绪二十三年（1897年），围绕「德国以巨野教案为由占领胶州湾，清廷面临新的租借压力」的照会、条约文本、换约或地方执行报告分别进入外交文书链。核心取证：《清德宗实录》光绪二十三年相关条；《筹办夷务始末》相关年卷与中外照会、条约文本。这一锚点记录的是文书链：它能证明呈报、上谕、部议或照会进入机构流程，不能由此推出实际执行。来源保留：此行仅确认相关文书链和可追查的行政动作；不得由其直接推断政策有效、地方服从、民众态度或单一因果。
\eventanchor{EQ-1898-LEASED-TERRITORIES-DOCUMENTS}{围绕「俄、英、法、德相继取得租…}光绪二十四年（1898年），围绕「俄、英、法、德相继取得租借地或势力范围安排」的照会、条约文本、换约或地方执行报告分别进入外交文书链。核心取证：《清德宗实录》光绪二十四年相关条；《筹办夷务始末》相关年卷与中外照会、条约文本。这一锚点记录的是文书链：它能证明呈报、上谕、部议或照会进入机构流程，不能由此推出实际执行。来源保留：此行仅确认相关文书链和可追查的行政动作；不得由其直接推断政策有效、地方服从、民众态度或单一因果。
\eventanchor{EQ-1898-HUNDRED-DAYS-REFORM-DOCUMENTS}{围绕「光绪帝颁布变法诏令，推动…}光绪二十四年（1898年），围绕「光绪帝颁布变法诏令，推动官制、教育、经济与军事改革」的上谕、奏折与部议在中央—地方行政链中流转。核心取证：《清德宗实录》光绪二十四年相关条；宫中朱批奏折、内阁大库题本与《清史稿》相关志传。这一锚点记录的是文书链：它能证明呈报、上谕、部议或照会进入机构流程，不能由此推出实际执行。来源保留：此行仅确认相关文书链和可追查的行政动作；不得由其直接推断政策有效、地方服从、民众态度或单一因果。
\eventanchor{EQ-1898-WUXU-COUP-DOCUMENTS}{围绕「慈禧发动政变，光绪帝被幽…}光绪二十四年（1898年），围绕「慈禧发动政变，光绪帝被幽禁，戊戌变法主要措施被撤销或改写」的上谕、奏折与部议在中央—地方行政链中流转。核心取证：《清德宗实录》光绪二十四年相关条；宫中朱批奏折、内阁大库题本与《清史稿》相关志传。这一锚点记录的是文书链：它能证明呈报、上谕、部议或照会进入机构流程，不能由此推出实际执行。来源保留：此行仅确认相关文书链和可追查的行政动作；不得由其直接推断政策有效、地方服从、民众态度或单一因果。
\eventanchor{EQ-1899-BOXER-UPRISING-DOCUMENTS}{围绕「义和团在华北扩展，反教、…}光绪二十五年（1899年），围绕「义和团在华北扩展，反教、反洋冲突加剧」的地方呈报、案件或赈务记录将社会反应带入官府文书链。核心取证：《清德宗实录》光绪二十五年相关条；地方志、督抚奏折及地方档案。这一锚点记录的是文书链：它能证明呈报、上谕、部议或照会进入机构流程，不能由此推出实际执行。来源保留：此行仅确认相关文书链和可追查的行政动作；不得由其直接推断政策有效、地方服从、民众态度或单一因果。
\subsection{1900—1909}
\eventanchor{EQ-1900-BOXER-WAR-DOCUMENTS}{清廷围绕「慈禧对列强宣战，八国…}光绪二十六年（1900年），清廷围绕「慈禧对列强宣战，八国联军入侵并占领北京」调遣军队、筹措饷械并处理战报，中央与地方的军政文书形成连续记录。核心取证：《清德宗实录》光绪二十六年相关条；军机处档、战事方略及地方奏报。这一锚点记录的是文书链：它能证明呈报、上谕、部议或照会进入机构流程，不能由此推出实际执行。来源保留：此行仅确认相关文书链和可追查的行政动作；不得由其直接推断政策有效、地方服从、民众态度或单一因果。
\eventanchor{EQ-1900-SOUTHEAST-MUTUAL-PROTECTION-DOCUMENTS}{围绕「东南互保使部分督抚与列强…}光绪二十六年（1900年），围绕「东南互保使部分督抚与列强协商维持地方秩序」的上谕、奏折与部议在中央—地方行政链中流转。核心取证：《清德宗实录》光绪二十六年相关条；宫中朱批奏折、内阁大库题本与《清史稿》相关志传。这一锚点记录的是文书链：它能证明呈报、上谕、部议或照会进入机构流程，不能由此推出实际执行。来源保留：此行仅确认相关文书链和可追查的行政动作；不得由其直接推断政策有效、地方服从、民众态度或单一因果。
\eventanchor{EQ-1901-BOXER-PROTOCOL-DOCUMENTS}{围绕「《辛丑条约》签订，赔款、…}光绪二十七年（1901年），围绕「《辛丑条约》签订，赔款、驻军和外交条件进一步限制清廷」的照会、条约文本、换约或地方执行报告分别进入外交文书链。核心取证：《清德宗实录》光绪二十七年相关条；《筹办夷务始末》相关年卷与中外照会、条约文本。这一锚点记录的是文书链：它能证明呈报、上谕、部议或照会进入机构流程，不能由此推出实际执行。来源保留：此行仅确认相关文书链和可追查的行政动作；不得由其直接推断政策有效、地方服从、民众态度或单一因果。
\eventanchor{EQ-1901-ZONGLI-YAMEN-REPLACED-DOCUMENTS}{围绕「总理衙门改为外务部，晚清…}光绪二十七年（1901年），围绕「总理衙门改为外务部，晚清外交机构进一步制度化」的上谕、奏折与部议在中央—地方行政链中流转。核心取证：《清德宗实录》光绪二十七年相关条；宫中朱批奏折、内阁大库题本与《清史稿》相关志传。这一锚点记录的是文书链：它能证明呈报、上谕、部议或照会进入机构流程，不能由此推出实际执行。来源保留：此行仅确认相关文书链和可追查的行政动作；不得由其直接推断政策有效、地方服从、民众态度或单一因果。
\eventanchor{EQ-1902-CIXI-RETURN-BEIJING-DOCUMENTS}{围绕「慈禧与光绪帝返回北京，战…}光绪二十八年（1902年），围绕「慈禧与光绪帝返回北京，战后朝廷开始推出新政措施」的上谕、奏折与部议在中央—地方行政链中流转。核心取证：《清德宗实录》光绪二十八年相关条；宫中朱批奏折、内阁大库题本与《清史稿》相关志传。这一锚点记录的是文书链：它能证明呈报、上谕、部议或照会进入机构流程，不能由此推出实际执行。来源保留：此行仅确认相关文书链和可追查的行政动作；不得由其直接推断政策有效、地方服从、民众态度或单一因果。
\eventanchor{EQ-1902-NEW-POLICIES-EDUCATION-DOCUMENTS}{围绕「清末新政在学制、军制、商…}光绪二十八年（1902年），围绕「清末新政在学制、军制、商政和地方行政方面逐步启动」的禁令、编纂、教育或传播活动留下中央和地方文本记录。核心取证：《清德宗实录》光绪二十八年相关条；内阁档、文集或报刊相关记录。这一锚点记录的是文书链：它能证明呈报、上谕、部议或照会进入机构流程，不能由此推出实际执行。来源保留：此行仅确认相关文书链和可追查的行政动作；不得由其直接推断政策有效、地方服从、民众态度或单一因果。
\eventanchor{EQ-1903-BANNER-POSTS-REFORM-DOCUMENTS}{围绕「清廷废除部分旗人官职垄断…}光绪二十九年（1903年），围绕「清廷废除部分旗人官职垄断，调整旗籍制度与职务资格」的地方呈报、案件或赈务记录将社会反应带入官府文书链。核心取证：《清德宗实录》光绪二十九年相关条；地方志、督抚奏折及地方档案。这一锚点记录的是文书链：它能证明呈报、上谕、部议或照会进入机构流程，不能由此推出实际执行。来源保留：此行仅确认相关文书链和可追查的行政动作；不得由其直接推断政策有效、地方服从、民众态度或单一因果。
\eventanchor{EQ-1903-SHANGHAI-COMMERCIAL-LAW-DOCUMENTS}{围绕「商法、公司与工商管理的近…}光绪二十九年（1903年），围绕「商法、公司与工商管理的近代制度讨论在口岸和中央之间推进」的经费、税收、赈济或工程呈报经由户部与地方衙门处理。核心取证：《清德宗实录》光绪二十九年相关条；户部奏销、盐政、河工或海关档案。这一锚点记录的是文书链：它能证明呈报、上谕、部议或照会进入机构流程，不能由此推出实际执行。来源保留：此行仅确认相关文书链和可追查的行政动作；不得由其直接推断政策有效、地方服从、民众态度或单一因果。
\eventanchor{EQ-1904-RUSSO-JAPANESE-WAR-DOCUMENTS}{清廷围绕「日俄战争在中国东北爆…}光绪三十年（1904年），清廷围绕「日俄战争在中国东北爆发，清廷名义中立而东北遭受战场化」调遣军队、筹措饷械并处理战报，中央与地方的军政文书形成连续记录。核心取证：《清德宗实录》光绪三十年相关条；军机处档、战事方略及地方奏报。这一锚点记录的是文书链：它能证明呈报、上谕、部议或照会进入机构流程，不能由此推出实际执行。来源保留：此行仅确认相关文书链和可追查的行政动作；不得由其直接推断政策有效、地方服从、民众态度或单一因果。
\eventanchor{EQ-1904-IMPERIAL-EXAMINATION-DEBATE-DOCUMENTS}{围绕「科举存废与新式教育的争论…}光绪三十年（1904年），围绕「科举存废与新式教育的争论加速」的禁令、编纂、教育或传播活动留下中央和地方文本记录。核心取证：《清德宗实录》光绪三十年相关条；内阁档、文集或报刊相关记录。这一锚点记录的是文书链：它能证明呈报、上谕、部议或照会进入机构流程，不能由此推出实际执行。来源保留：此行仅确认相关文书链和可追查的行政动作；不得由其直接推断政策有效、地方服从、民众态度或单一因果。
\eventanchor{EQ-1905-ABOLITION-OF-CIVIL-SERVICE-EXAM-DOCUMENTS}{围绕「清廷废除科举，延续千余年…}光绪三十一年（1905年），围绕「清廷废除科举，延续千余年的取士制度终结」的禁令、编纂、教育或传播活动留下中央和地方文本记录。核心取证：《清德宗实录》光绪三十一年相关条；内阁档、文集或报刊相关记录。这一锚点记录的是文书链：它能证明呈报、上谕、部议或照会进入机构流程，不能由此推出实际执行。来源保留：此行仅确认相关文书链和可追查的行政动作；不得由其直接推断政策有效、地方服从、民众态度或单一因果。
\eventanchor{EQ-1905-TONGMENGHUI-DOCUMENTS}{围绕「中国同盟会在东京成立，革…}光绪三十一年（1905年），围绕「中国同盟会在东京成立，革命组织和宣传网络进一步整合」的地方呈报、案件或赈务记录将社会反应带入官府文书链。核心取证：《清德宗实录》光绪三十一年相关条；地方志、督抚奏折及地方档案。这一锚点记录的是文书链：它能证明呈报、上谕、部议或照会进入机构流程，不能由此推出实际执行。来源保留：此行仅确认相关文书链和可追查的行政动作；不得由其直接推断政策有效、地方服从、民众态度或单一因果。
\eventanchor{EQ-1906-CONSTITUTIONAL-PROMISE-DOCUMENTS}{围绕「清廷宣布预备立宪，着手考…}光绪三十二年（1906年），围绕「清廷宣布预备立宪，着手考察各国制度和改革官制」的上谕、奏折与部议在中央—地方行政链中流转。核心取证：《清德宗实录》光绪三十二年相关条；宫中朱批奏折、内阁大库题本与《清史稿》相关志传。这一锚点记录的是文书链：它能证明呈报、上谕、部议或照会进入机构流程，不能由此推出实际执行。来源保留：此行仅确认相关文书链和可追查的行政动作；不得由其直接推断政策有效、地方服从、民众态度或单一因果。
\eventanchor{EQ-1906-NEW-ARMY-REFORM-DOCUMENTS}{清廷围绕「新军编练与军事学堂扩…}光绪三十二年（1906年），清廷围绕「新军编练与军事学堂扩展，军队成为新式政治动员的重要空间」调遣军队、筹措饷械并处理战报，中央与地方的军政文书形成连续记录。核心取证：《清德宗实录》光绪三十二年相关条；军机处档、战事方略及地方奏报。这一锚点记录的是文书链：它能证明呈报、上谕、部议或照会进入机构流程，不能由此推出实际执行。来源保留：此行仅确认相关文书链和可追查的行政动作；不得由其直接推断政策有效、地方服从、民众态度或单一因果。
\eventanchor{EQ-1907-MANCHURIAN-PROVINCES-DOCUMENTS}{清廷围绕「清廷将东三省调整为行…}光绪三十三年（1907年），清廷围绕「清廷将东三省调整为行省体制，加强近代行政与边防经营」处理盟旗、驻防、交通或地方呈报，相关边务文书进入中央决策链。核心取证：《清德宗实录》光绪三十三年相关条；军机处档、理藩院档或相关方略。这一锚点记录的是文书链：它能证明呈报、上谕、部议或照会进入机构流程，不能由此推出实际执行。来源保留：此行仅确认相关文书链和可追查的行政动作；不得由其直接推断政策有效、地方服从、民众态度或单一因果。
\eventanchor{EQ-1907-QING-LEGAL-REFORM-DOCUMENTS}{围绕「清廷推进修律、审判和警政…}光绪三十三年（1907年），围绕「清廷推进修律、审判和警政改革，尝试改造传统法律制度」的上谕、奏折与部议在中央—地方行政链中流转。核心取证：《清德宗实录》光绪三十三年相关条；宫中朱批奏折、内阁大库题本与《清史稿》相关志传。这一锚点记录的是文书链：它能证明呈报、上谕、部议或照会进入机构流程，不能由此推出实际执行。来源保留：此行仅确认相关文书链和可追查的行政动作；不得由其直接推断政策有效、地方服从、民众态度或单一因果。
\eventanchor{EQ-1908-GUANGXU-CIXI-DEATHS-DOCUMENTS}{围绕「光绪帝与慈禧太后相继去世…}光绪三十四年（1908年），围绕「光绪帝与慈禧太后相继去世，溥仪继位为宣统帝」的上谕、奏折与部议在中央—地方行政链中流转。核心取证：《清德宗实录》光绪三十四年相关条；宫中朱批奏折、内阁大库题本与《清史稿》相关志传。这一锚点记录的是文书链：它能证明呈报、上谕、部议或照会进入机构流程，不能由此推出实际执行。来源保留：此行仅确认相关文书链和可追查的行政动作；不得由其直接推断政策有效、地方服从、民众态度或单一因果。
\eventanchor{EQ-1908-CONSTITUTIONAL-OUTLINE-DOCUMENTS}{围绕「清廷公布《钦定宪法大纲》…}光绪三十四年（1908年），围绕「清廷公布《钦定宪法大纲》等文件，明确预备立宪框架」的上谕、奏折与部议在中央—地方行政链中流转。核心取证：《清德宗实录》光绪三十四年相关条；宫中朱批奏折、内阁大库题本与《清史稿》相关志传。这一锚点记录的是文书链：它能证明呈报、上谕、部议或照会进入机构流程，不能由此推出实际执行。来源保留：此行仅确认相关文书链和可追查的行政动作；不得由其直接推断政策有效、地方服从、民众态度或单一因果。
